% =====================================================================
%  Chapter 9 — Four Behavioural Traps and Their Magnitudes
% =====================================================================

\chapter{Four Behavioural Traps and Their Magnitudes}
\label{ch:behaviour}

Chapter 8's effect is about how often the investor looks. This chapter is
about what the investor does when they do. The four traps below are the
ones the literature has quantified, and each one has a number attached,
because a trap without a magnitude is an anecdote.

The literature these traps come from is broader than the four papers
below. Shiller's \emph{Irrational Exuberance} is the standard account of
the underlying phenomenon: that prices are not the output of a rational
expectations process, that narratives move them, and that the same
mechanisms that produce exuberance produce the errors an individual
investor is about to make. The four traps are the individual-level,
quantified pieces of it.

\section{Trap one: the disposition effect}

Odean examined 10,000 brokerage accounts from 1987 to 1993 and found that
investors realise gains far more readily than losses: the average annual
return on stocks sold was 14.83 per cent, against 9.76 per cent on stocks
held. And the effect persists: in the year after a sale, the winners that
had been sold outperformed the losers that had been held by 3.4 percentage
points.

That 3.4-point figure is the size of the opportunity that the trap denies
the investor. It is not a small effect and it is not a rounding error; it
is comparable to the entire equity premium for the decade.

The mechanism is now modelled. Barberis and Xiong's 2009 analysis in the
\emph{Journal of Finance} shows that a preference for realising gains
combined with a loss-aversion term that depends on the investor's own
realised gains and losses generates the disposition effect as a consequence,
rather than treating it as a primitive. The model also reproduces the two
facts that make the effect economically consequential: realised losses are
small in absolute terms relative to the portfolio, so the tax cost of
realising them is low; and the portfolio's position in the losing asset is
small, so the expected benefit from realising the loss is also small. The
gains are being harvested for the wrong reason and the losses are being
suffered for the wrong reason, and the investor is worse off for it.

The implication for a system is specific: \emph{never let a rule depend on
whether a position is currently above or below its cost basis}. The tax
engine in Chapter 29 tracks lots and holding periods because the tax code
requires it, but no allocation rule in this book is permitted to read the
cost basis. A rule that could see the cost basis would be able to reproduce
the disposition effect exactly, and would have to be argued out of it.

\section{Trap two: trading too much}

Barber and Odean examined 66,465 households from 1991 to 1996 and found
that the most active traders --- the top quartile by turnover --- earned
11.4 per cent a year, against 17.9 per cent a year for the market. The
most active traders underperformed the market by 6.5 percentage points a
year, and the more active the household, the worse it did.

This is the trap that a machine-learning system is most obviously exposed
to, because a system that is permitted to rebalance without friction will
rebalance without reason. It is also the trap that a system with a
no-trade band, a turnover report and an explicit spread-cost model is
partially protected against --- and the protection is partial, not
complete, because the discipline has to be maintained by whoever runs it.

\section{Trap three: herding and the information cascade}

The third trap is social rather than cognitive, and it is the one the
least literature has quantified at the individual level. It operates
through the information cascade: when investors cannot distinguish a
private signal from others' actions, they follow, and the following is
itself the information that the next investor follows. The consequences are
large in the aggregate --- the herding observed in equity markets
contributes to the excess volatility documented in Chapter 3 --- and
individually they are the mechanism by which an investor's allocation
becomes a copy of a benchmark's, with all the benchmark's risks and none of
its scale.

For a single investor reading this book, the practical form of the trap is
the most obvious one: the temptation to hold what is fashionable, and the
tendency to believe that the reason it is fashionable is that it is good.
The structural defence is that the strategic allocation in this book is
\emph{stated in advance}, with a written justification, and is not
re-derived from whatever happened to work last year. That is not
intellectual timidity; it is the direct consequence of the evidence in
Chapter 3 that a policy fitted to the recent past is a policy fitted to
noise.

\section{Trap four: the illusion of control through sophistication}

The fourth trap is the one this book is most exposed to, and it deserves
the most space.

The temptation runs like this: the investor has a machine. The machine has
a model with twelve features and a Bayesian posterior and a Kelly sizing
rule. The outputs are precise, well-formatted, and update every month.
Precision of presentation is easily mistaken for precision of knowledge,
and a system that reports a Sharpe ratio to three decimal places invites
the reader to treat it as a measurement.

Two protections are built into this book, and both are visible in its
structure.

The first is that the selection machinery of Part IV is applied to this
book's own family of configurations and its own results are reported
whatever they are. The second is that the trial family is declared in
advance and its size is published, so that a reader can compute for
themselves how much of the reported performance is the maximum of a noisy
sample rather than a property of the system.

\begin{takeaway}[The four traps and their sizes]
\begin{itemize}
\item \textbf{Disposition effect:} winners sold beat losers held by 3.4
percentage points in the following year.
\item \textbf{Overtrading:} the most active quartile of retail households
earned 11.4 per cent against the market's 17.9 per cent.
\item \textbf{Herding:} the aggregate mechanism is well documented; the
individual defence is a pre-commitment to a stated allocation.
\item \textbf{Sophistication:} the defence is a declared trial family and
the willingness to publish a null result.
\end{itemize}
\end{takeaway}

\section{Why a system is not immune to any of them}

It is worth being explicit that a machine-learning system is not immune to
these traps, and in two of the four it is more exposed than a human.

The disposition effect cannot bite a system directly, because the system
does not experience a gain or a loss. But it can bite the system's
\emph{owner}: if the system holds a losing position and the owner sees the
account statement, the owner may override it. This is why the system's
reporting is designed to be read quarterly at most, and why the drawdown
brake exists --- to give the owner a reason not to override.

Overtrading is worse for a system than for a human, because a system that
is permitted to trade whenever its signal moves will trade whenever its
signal moves, and the signal moves constantly. The no-trade band and the
turnover statistic in Chapter 33 are the defences, and they are stated as
numbers so that the owner can see them.

Hereding does not apply to the system at all, and that is one of the
genuine advantages of automation: the system holds the stated allocation
because the allocation is stated, not because the last three months were
good.

Sophistication is the trap the system walks into by construction, and the
only defence is the one this book spends Part IV on.

\begin{keytakeaways}
\item The four traps have magnitudes: 3.4 points a year denied by the
disposition effect, 6.5 points a year lost by the most active retail
traders, plus a documented aggregate herding effect and the book's own
version of the sophistication trap.
\item A rule that could see the cost basis could reproduce the disposition
effect exactly; no allocation rule in this book is permitted to read it.
\item A system is more exposed than a human to overtrading and to the
illusion of control, and only partly protected by a no-trade band and a
declared trial family.
\item The defences are stated as numbers --- turnover, trial count,
published results --- so that a reader can check them rather than trust
them.
\end{keytakeaways}
